\subsection{三角関数の積分とのつながり}

$x=\sin^2t$と置くと，$0<t<\frac{\pi}{2}$で$\dd x=2\sin t\cos t\dd t$である。$\sin t,\cos t$が正なので，実数の指数もそのまま扱える。式\eqref{式:ベータ関数の定義}は
\begin{equation}\label{式:ベータ三角表示}
 B(p,q)=2\int_0^{\frac{\pi}{2}}
 (\sin t)^{2p-1}(\cos t)^{2q-1}\dd t
\end{equation}
となる。厳密には端点を避けて置換した後で極限を取る。左辺の両端の積分が収束しているので，対応する右辺の極限も有限である。

$p=q=\frac12$なら両方の指数が$0$になり，
\[
 B\left(\frac12,\frac12\right)
 =2\int_0^{\frac{\pi}{2}}1\dd t=\pi.
\]
つまり
\[
 \int_0^1\frac1{\sqrt{x(1-x)}}\dd x=\pi
\]
である。両端で高さが限りなく大きくなるにもかかわらず，積分は有限になる具体例である。

一方，$p=\frac{n+1}{2}$，$q=\frac12$とすれば
\[
 B\left(\frac{n+1}{2},\frac12\right)=2I_n.
\]
ベータ関数の漸化式で$p$を$1$増やすことは，$\sin$の指数$n$を$2$増やすことに対応する。実際
\[
 B\left(\frac{n+3}{2},\frac12\right)
 =\frac{n+1}{n+2}
  B\left(\frac{n+1}{2},\frac12\right)
\]
は$I_{n+2}=\frac{n+1}{n+2}I_n$そのものである。整数の多項式積分とウォリスの積分は，ベータ積分の異なる場合として結び付いた。

